\documentclass[11pt]{exam}
\usepackage[a4paper, margin=0.5in, footskip=0.1in]{geometry}
\usepackage{graphicx}
\usepackage{amsmath}
\usepackage{tabularx}
\usepackage{array}
\usepackage[version=4]{mhchem}
\usepackage{enumitem}
\usepackage{caption}
\usepackage{siunitx}
\usepackage{setspace}
\usepackage{tabto}
\usepackage{hyperref}
\usepackage{ragged2e}
\setlength{\parindent}{0pt}

\footer{}{\thepage}{}
\renewcommand{\questionlabel}{\textbf{\thequestion}\hfill}
\renewcommand{\partlabel}{\textbf{(\thepartno)}\hfill}
\renewcommand{\subpartlabel}{\textbf{(\thesubpart)}}
\renewcommand{\choicelabel}{\textbf{\thechoice}\hspace{0.5em}}
\renewcommand{\thesubsubpart}{\arabic{subsubpart}}
\renewcommand{\subsubpartlabel}{\textbf{\thesubsubpart.}}
\makeatletter
\renewcommand{\dotfill}{%
  \leavevmode\cleaders\hb@xt@ 3pt{\hss.\hss}\hfill\kern\z@
}

\renewcommand{\choiceshook}{%
  \setlength{\leftmargin}{1.7em}
}
\renewcommand{\questionshook}{
  \setlength{\labelsep}{0.2cm}
  \setlength{\leftmargin}{0.7cm}
}
\renewcommand{\partshook}{
\setlength{\labelsep}{0.35 cm}
\setlength{\leftmargin}{2em}
}
\renewcommand{\subpartshook}{
\setlength{\labelsep}{0.35 cm}
\setlength{\leftmargin}{1.8em}
}
\renewcommand{\subsubpartshook}{
\setlength{\labelsep}{0.35 cm}
\setlength{\leftmargin}{1.6em}
}

\setlength{\parindent}{0pt}

\sisetup{detect-all = true}
\renewcommand{\choiceshook}{%
  \setlength{\leftmargin}{1.7em}
}
\renewcommand{\questionshook}{
  \setlength{\labelsep}{0.5em}
  \setlength{\leftmargin}{0.8cm}
}
\renewcommand{\partshook}{
\setlength{\labelsep}{1em}
\setlength{\leftmargin}{1.7em}
}

\title{2018 NY Catholic High School Year 4 Chemistry EOY}
\date{}

\begin{document}
\singlespacing
\setcounter{page}{1}

\begin{center}
{\Large \textbf{NY CATHOLIC HIGH SCHOOL}} \\
{\Large \textbf{Year 4 Integrated Programme}} \\
\vspace{1em}
{\Large \textbf{END-OF-YEAR EXAMINATION}} \\
{\Large \textbf{CHEMISTRY}} \\
\vspace{1em}
{\large Paper 1 Multiple Choice} \\
{\large 8 October 2018} \\
{\large 1 hour} \\
\end{center}

\vspace{1em}
{\bf Additional Materials:} Multiple Choice Answer Sheet

\vspace{0.5em}
{\bf READ THESE INSTRUCTIONS FIRST}

Write in soft pencil.

Do not use staples, paper clips, glue or correction fluid.

Write your name, class and index number on the Answer Sheet in the spaces provided.

There are forty questions in this paper. Answer all questions. For each question, there are four possible answers A, B, C and D.

Choose the one you consider correct and record your choice in soft pencil on the separate Answer Sheet.

Read the instruction on the Answer Sheet very carefully.

Each correct answer will score one mark. A mark will not be deducted for a wrong answer.

Any rough working should be done in this booklet.

A copy of the Periodic Table is printed on page~16.

The use of an approved scientific calculator is expected, where appropriate.

\newpage

\section*{Paper 1}

\begin{questions}

\question Gas Y is less dense than air, very soluble in water and is an alkali. Which method is used to collect a dry sample of the gas?

\textbf{[Diagram omitted]}

\begin{choices}
\choice A -- dry Y collected over water with calcium oxide
\choice B -- dry Y collected over water with calcium oxide
\choice C -- Apparatus with porous pot showing diffusion of gases
\choice D -- dry Y collected over concentrated sulfuric acid
\end{choices}

\question The apparatus below is used to show the diffusion of gases. \\
\textbf{[Diagram omitted]}

Which pair of gases X and Y would cause no movement of the water level at Z?

\begin{tabular}{|c|c|c|}
\hline
 & X & Y \\
\hline
A & \ce{CO} & \ce{C2H6} \\
B & \ce{CO} & \ce{CO2} \\
C & \ce{C2H4} & \ce{C2H6} \\
D & \ce{NO} & \ce{C2H4} \\
\hline
\end{tabular}

\question Use the information to answer questions 3 and 4.

The filter tip of a cigarette acts as both a filter and a condenser. You may assume the filter tip is 100\% efficient.

\begin{tabular}{|c|c|c|}
\hline
 & substance & boiling point/\qty{}{\celsius} \\
\hline
1 & carbon monoxide & -191 \\
2 & nicotine & 247 \\
3 & tar & 350 to 400 \\
4 & water & 100 \\
\hline
\end{tabular}

Which substance cannot be removed?

\begin{choices}
\choice 1 only
\choice 1 and 4 only
\choice 2 and 3 only
\choice 2, 3 and 4 only
\end{choices}

\question With reference to the table provided in question 3, which is an impure substance?

\begin{choices}
\choice carbon monoxide and nicotine only
\choice carbon monoxide, nicotine and tar only
\choice carbon monoxide, nicotine and water only
\choice tar only
\end{choices}

\newpage

\question A sample of a molten pure compound was cooled. The graph shows how the temperature of the compound changes with time as it was cooled.

\textbf{[Diagram omitted]}

At which region(s) were all the particles of the compound vibrating about a fixed position?

\begin{choices}
\choice P to Q
\choice Q to R
\choice R to S
\choice Q to R and R to S
\end{choices}

\question Which substance is a compound that exists as molecules?

\begin{choices}
\choice aluminium oxide
\choice brass
\choice oxygen
\choice sulfur dioxide
\end{choices}

\question Some isotopes are unstable and decompose naturally to form another element. In one type of decomposition, a neutron in the nucleus decomposes to form two sub-atomic particles, a proton, which is retained in the nucleus, and an electron, which is expelled from the atom.

Which change describes this process?

\begin{choices}
\choice \ce{^{11}_{6}C -> ^{12}_{7}N}
\choice \ce{^{22}_{11}Na -> ^{22}_{10}Ne}
\choice \ce{^{32}_{15}P -> ^{32}_{16}P}
\choice \ce{^{40}_{20}Ca -> ^{40}_{20}Ca}
\end{choices}

\question The formulae of the ions of four elements found in period 2 of the Periodic Table are shown below.

W\textsuperscript{2-}\quad X\textsuperscript{-}\quad Y\textsuperscript{+}\quad Z\textsuperscript{2+}

Which statement about these ions is incorrect?

\begin{choices}
\choice Element W has a low melting and boiling point.
\choice Element Y is soft and malleable.
\choice X is an oxidising agent.
\choice X and Z form an ionic compound \ce{ZX2}.
\end{choices}

\question The melting point of calcium oxide differs from that of sodium chloride. Which statement best explains this?

\begin{choices}
\choice A calcium ion has more protons than a sodium ion.
\choice Calcium oxide has a lower relative molecular mass than sodium chloride.
\choice Ions in calcium oxide have stronger electrostatic forces of attraction than sodium chloride.
\choice Sodium is more reactive than calcium.
\end{choices}

\question Which property is \textbf{not} typical of transition metals?

\begin{choices}
\choice They form coloured compounds.
\choice They have high melting points.
\choice They have low densities.
\choice They show variable oxidation states.
\end{choices}

\newpage

\question A dative covalent bond is formed when one atom provides two electrons for sharing. Which molecule contains a dative covalent bond?

\begin{choices}
\choice ammonia
\choice carbon monoxide
\choice hydrogen chloride
\choice sulfur dioxide
\end{choices}

\question The physical properties of four substances are shown in the table below.

\begin{tabular}{|c|c|c|c|c|c|}
\hline
substance & melting point/\qty{}{\celsius} & boiling point/\qty{}{\celsius} & density/\qty{}{\g\per\cm\cubed} & \multicolumn{2}{c|}{electrical conductivity} \\
\cline{5-6}
 &  &  &  & in solid state & in liquid state \\
\hline
P & -219 & -183 & 1.14 & \texttimes & \texttimes \\
Q & 98 & 883 & 0.968 & \checkmark & \checkmark \\
R & 801 & 1413 & 2.17 & \texttimes & \checkmark \\
S & 0 & 100 & 1.00 & \texttimes & \texttimes \\
\hline
\end{tabular}

\vspace{0.5cm}
legend: \texttimes{} = cannot conduct electricity \quad \checkmark{} = can conduct electricity

Which statement(s) is/are correct?

1. P has a simple molecular structure. \\
2. Q contains mobile ions in the solid state. \\
3. R has a giant ionic crystal lattice structure. \\
4. S exhibits hydrogen bonding with propanol molecules.

\begin{choices}
\choice All of the above
\choice 1, 2 and 3 only
\choice 1, 3 and 4 only
\choice 1 and 3 only
\end{choices}

\question The pollutant sulfur dioxide, \ce{SO2}, can be removed in desulfurisation plants with the use of chemicals such as calcium carbonate or magnesium hydroxide. These chemicals react in a similar process known as scrubbing to remove the \ce{SO2}.

What is the main product formed initially when magnesium hydroxide reacts with \ce{SO2}?

\begin{choices}
\choice \ce{Mg(HSO3)2}
\choice \ce{MgS}
\choice \ce{MgSO3}
\choice \ce{MgSO4}
\end{choices}

\question Which step in the diagram shows the process of photosynthesis?

\textbf{[Diagram omitted: carbon cycle showing carbon dioxide in atmosphere, carbon compounds in plants (labelled B), carbon compounds in animals (labelled C), and oil and natural gas (labelled D), with arrow A from atmosphere to plants]}

\begin{choices}
\choice A
\choice B
\choice C
\choice D
\end{choices}

\newpage

\question Element R reacts with oxygen to form a gas, T. T changes the colour of damp litmus paper from blue to red. T is used to kill bacteria in the preservation of dried fruit. T can be neutralised by calcium oxide.

What is R?

\begin{choices}
\choice carbon
\choice chlorine
\choice nitrogen
\choice sulfur
\end{choices}

\question Element X sinks in water. It reacts with cold water and steam. The oxide of element X cannot be reduced by carbon. The carbonate of element X takes a longer time than \ce{ZnCO3} to thermally decompose.

Which details are those of X?

\begin{tabular}{|c|c|c|}
\hline
 & proton number & mass number \\
\hline
A & 11 & 23 \\
B & 20 & 40 \\
C & 29 & 64 \\
D & 82 & 207 \\
\hline
\end{tabular}

\question The reaction of substance X is shown.

\begin{tabular}{|l|l|}
\hline
Reaction & Observation \\
\hline
X is placed in 2~cm\textsuperscript{3} dilute hydrochloric acid in a test tube. & Effervescence of a colourless odourless gas. \\
\hline
A lighted splint is immediately inserted into the test tube. & Lighted splint extinguishes. No `pop' sound was heard. \\
\hline
\end{tabular}

Which substance could X be?

\begin{choices}
\choice aqueous ammonium nitrate
\choice silicon(IV) oxide
\choice sodium carbonate
\choice zinc
\end{choices}

\question Pyridine, \ce{C5H5N}, is an organic compound which is a weak base. It has similar chemical properties to aqueous ammonia.

Which reagent will show no observable change with pyridine?

\begin{choices}
\choice aqueous aluminium nitrate
\choice aqueous calcium nitrate
\choice aqueous copper(II) sulfate
\choice aqueous zinc chloride
\end{choices}

\question Which substance, in \qty{1.0}{\mol\per\dm\cubed} aqueous solution, would have the same hydrogen ion concentration as \qty{2.0}{\mol\per\dm\cubed} of hydrochloric acid?

\begin{choices}
\choice ethanoic acid
\choice nitric acid
\choice phosphoric acid
\choice sulfuric acid
\end{choices}

\newpage

\question Most salts may be prepared by reacting a metal, base or a carbonate with a dilute acid.

Which salt should \textbf{not} be prepared with the following reactants?

\begin{tabular}{|c|c|c|c|}
\hline
 & reactant 1 & reactant 2 & salt to be prepared \\
\hline
A & \ce{BaCO3} & \ce{H2SO4} & \ce{BaSO4} \\
B & \ce{CuO} & \ce{H2SO4} & \ce{CuSO4} \\
C & \ce{NaOH} & \ce{HNO3} & \ce{NaNO3} \\
D & \ce{Zn} & \ce{HCl} & \ce{ZnCl2} \\
\hline
\end{tabular}

\question A student carried out the following series of tests on an aqueous solution of potassium carbonate and recorded his results in the table below.

Which test should be repeated because of an incorrect observation in the table?

\begin{tabular}{|c|c|c|}
\hline
 & test & observation \\
\hline
A & add barium chloride solution & white precipitate formed \\
B & add calcium nitrate solution & white precipitate formed \\
C & add dilute hydrochloric acid & effervescence \\
D & add sodium iodide solution & yellow precipitate formed \\
\hline
\end{tabular}

\question Four separate solutions are prepared so that each solution contains 1~g of the ions \ce{Fe^{2+}}, \ce{Fe^{3+}}, \ce{Ca^{2+}} and \ce{Zn^{2+}}. An excess of aqueous sodium hydroxide solution is added, and the mass of precipitate formed is determined.

\textbf{[Diagram omitted: four graphs A--D of mass of precipitate vs ion type]}

Which graph illustrates the results?

\begin{choices}
\choice A
\choice B
\choice C
\choice D
\end{choices}

\question In the contact process for making sulfuric acid, one step involves the oxidation of sulfur dioxide as shown.

\ce{2SO2(g) + O2(g) <=> 2SO3(g)} \qquad $\Delta H = \qty{-197}{\kJ\per\mol}$

Which change will increase the yield of sulfur trioxide?

\begin{choices}
\choice adding a catalyst
\choice decreasing the pressure
\choice decreasing the temperature
\choice removing \ce{SO2}
\end{choices}

\newpage

\question \qty{480}{\cm\cubed} of carbon dioxide was reacted with \qty{0.12}{\g} of carbon. Calculate the total volume of gases remaining after the reaction. All measurements are done at room temperature and pressure. [\ce{A_r}: C, 12; O, 16]

\ce{CO2(g) + C(s) -> 2CO(g)}

\begin{choices}
\choice \qty{240}{\cm\cubed}
\choice \qty{480}{\cm\cubed}
\choice \qty{720}{\cm\cubed}
\choice \qty{960}{\cm\cubed}
\end{choices}

\question Egg shell contains mainly calcium carbonate. \qty{5.0}{\g} of impure egg shell was reacted with an excess of dilute nitric acid. \qty{0.6}{\dm\cubed} of gas and \qty{0.45}{\g} of water was collected.

What is the percentage purity of the sample of the egg shell?

\begin{choices}
\choice 1.2\%
\choice 9.0\%
\choice 50\%
\choice 63\%
\end{choices}

\question Three experiments to investigate the reactivity of three metals are shown.

\textbf{[Diagram omitted: manganese with aqueous nickel(II) nitrate -- nickel displaced; chromium with aqueous nickel(II) nitrate -- nickel displaced; manganese with aqueous manganese(II) nitrate -- not displaced]}

What is the correct order of reactivity for these three metals?

\begin{choices}
\choice chromium > manganese > nickel
\choice manganese > chromium > nickel
\choice manganese > nickel > chromium
\choice nickel > chromium > manganese
\end{choices}

\question Hydrogen gas is allowed to react with copper(II) oxide. After the reaction is complete, the burner is turned off, but the flow of hydrogen gas is continued until the tube is cool.

\textbf{[Diagram omitted: apparatus showing hydrogen gas passing over heated copper(II) oxide]}

Why is the hydrogen gas allowed to flow through the tube during cooling?

\begin{choices}
\choice to allow the tube to cool slowly
\choice to prevent an explosion in the hot tube
\choice to prevent copper from reacting with the air
\choice to remove any traces of water left in the tube
\end{choices}

\newpage

\question The diagram shows a boat made from iron. Some magnesium blocks are attached to the iron below the water line.

\textbf{[Diagram omitted: boat with magnesium blocks attached to iron hull below water line]}

Why does the magnesium stop the iron from rusting?

\begin{choices}
\choice Magnesium loses electrons more readily than iron.
\choice Magnesium reacts to form a protective coating of magnesium oxide on the iron.
\choice The magnesium forms an alloy with the iron.
\choice The magnesium stops oxygen in the water from getting to the iron.
\end{choices}

\question Which equation represents a redox reaction?

\begin{choices}
\choice \ce{Ag+(aq) + Cl-(aq) -> AgCl(s)}
\choice \ce{ZnCO3(s) -> ZnO(s) + CO2(g)}
\choice \ce{H+(aq) + OH-(aq) -> H2O(l)}
\choice \ce{Mg(s) + 2H+(aq) -> Mg^{2+}(aq) + H2(g)}
\end{choices}

\question Four experiments on rusting are shown.

\textbf{[Diagram omitted: four test tubes -- 1: dry air + paper-clip, 15\celsius, not rusty; 2: tap water + paper-clip, 15\celsius, rusts; 3: boiled tap water + oil + paper-clip, 15\celsius, not rusty; 4: tap water + paper-clip, 25\celsius, rusts]}

Which two experiments can be used to show that air is needed for iron to rust?

\begin{choices}
\choice 1 and 3 only
\choice 1 and 4 only
\choice 2 and 3 only
\choice 2 and 4 only
\end{choices}

\question Which equation represents an endothermic reaction?

\begin{choices}
\choice \ce{CaCO3(s) -> CaO(s) + CO2(g)}
\choice \ce{2C2H6(g) + 7O2(g) -> 4CO2(g) + 6H2O(l)}
\choice \ce{C6H12O6(s) + 6O2(g) -> 6CO2(g) + 6H2O(g)}
\choice \ce{HCl(g) + NaOH(aq) -> H2O(l) + NaCl(aq)}
\end{choices}

\newpage

\question The diagram shows an apparatus used to electroplate a metal ring with copper.

\textbf{[Diagram omitted: electroplating apparatus with copper electrode, power supply, bulb, aqueous copper(II) sulfate electrolyte]}

The experiment did not work. What change is needed in the experiment to make it work?

\begin{choices}
\choice Add solid copper(II) sulfate to the electrolyte.
\choice Increase the temperature of the electrolyte.
\choice Replace the copper electrode with a carbon electrode.
\choice Reverse the connection of the battery.
\end{choices}

\question In which compound does the element Q have the highest oxidation state?

\begin{choices}
\choice \ce{QO}
\choice \ce{QO2}
\choice \ce{Q2O3}
\choice \ce{QO3}
\end{choices}

\question In an electrolysis experiment, the same amount of charge deposited \qty{24}{\g} of titanium and \qty{64}{\g} of copper. [\ce{A_r}: Ti, 48; Cu, 64]

What is the charge on the titanium ion?

\begin{choices}
\choice 1+
\choice 2+
\choice 3+
\choice 4+
\end{choices}

\question Aqueous copper(II) sulfate was electrolysed using copper electrodes.

Which observations were made?

\begin{tabular}{|c|c|c|c|}
\hline
 & at the anode & at the cathode & electrolyte \\
\hline
A & anode dissolves & pink-brown solid forms & blue solution fades \\
B & anode dissolves & pink-brown solid forms & no observable change \\
C & effervescence of colourless odourless gas & effervescence of colourless odourless gas & no observable change \\
D & effervescence of colourless odourless gas & pink-brown solid forms & blue solution fades \\
\hline
\end{tabular}

\newpage

\question The diagram shows the structure of the compound 1,3-butadiene (\ce{M_r} = \qty{54}{\g\per\mol}).

\textbf{[Diagram omitted: structure of \ce{H2C=CH-CH=CH2}]}

How many moles of hydrogen gas are needed to fully react with \qty{108}{\g} of 1,3-butadiene?

\begin{choices}
\choice 1
\choice 2
\choice 3
\choice 4
\end{choices}

\question Which compound exhibits cis-trans isomerism?

\textbf{[Diagram omitted: four compounds structures A, B, C, D]}

\begin{choices}
\choice A
\choice B
\choice C
\choice D
\end{choices}

\question The structure of a monomer is shown below:

\textbf{[Diagram omitted: monomer structure \ce{CH2=CH(CH3)Cl}]}

Which formula shows two repeating units of the polymer made from this monomer?

\textbf{[Diagram omitted: four polymer structures A--D]}

\begin{choices}
\choice A
\choice B
\choice C
\choice D
\end{choices}

\question The diagram represents a fat molecule. When fats undergo hydrolysis, their ester linkages are broken to form smaller molecules.

\textbf{[Diagram omitted: fat molecule structure]}

Which statement about the fat molecule shown is false?

\begin{choices}
\choice A longer alkyl chain on the fat molecule makes it less soluble in water.
\choice Addition of dilute hydrochloric acid will produce glycerol and fatty acid chains.
\choice It can undergo hydrogenation to produce margarine.
\choice It has a higher boiling point than a similar fat molecule that is completely saturated.
\end{choices}

\question The soap with the formula \ce{C17H35COONa} is

\begin{choices}
\choice a salt.
\choice an ester.
\choice formed from an unsaturated acid.
\choice formed from the hydrolysis of proteins.
\end{choices}

\end{questions}

\vfill
\begin{center}
\textbf{-- End of Paper 1 --}
\end{center}

\newpage

% ======================================================================
% SECTION A
% ======================================================================

\begin{center}
{\Large \textbf{CATHOLIC HIGH SCHOOL}} \\
\vspace{0.5em}
{\Large \textbf{End-of-Year Examination}} \\
{\Large \textbf{Year 4 (Integrated Programme)}} \\
{\Large \textbf{CHEMISTRY}} \\
\vspace{0.5em}
{\large 4 October 2018} \\
{\large 1 hour 45 minutes} \\
\end{center}

\vspace{0.5em}
Candidates answer on the Question Paper. \\
No Additional Materials are required.

\vspace{0.5em}
{\bf READ THESE INSTRUCTIONS FIRST}

Write your name, index number and class on all the work you hand in.

Write in dark blue or black pen.

You may use an HB pencil for any diagrams or graphs.

Do not use staples, paper clips, glue or correction fluid.

\vspace{1em}

\section*{Section A}

Answer \textbf{all} questions in the spaces provided.

\vspace{1em}

\begin{questions}

% ======================================================================
\question A1
\vspace{0.5em}

Choose from the following polymers to answer the questions.

\textbf{[Diagram omitted: structures of polymers A, B, C, D, E, F]}

Use the letters A to F to answer the following questions. Each polymer can be used once, more than once or not at all.

\begin{parts}

\part[1] Which polymer is used to make both clingfilm and plastic bags?

\vspace{1em}
\dotfill [1]

\part[2]
\begin{subparts}
\subpart[1] Give the letter of an addition polymer which is formed from propene.
\vspace{1em}
\dotfill [1]

\subpart[1] Give the letter of a polyester.
\vspace{1em}
\dotfill [1]
\end{subparts}

\part[1] Which polymer could be part of a protein?
\vspace{1em}
\dotfill [1]

\part[1] A common polymer used for home meal replacement trays is crystallised polyethylene terephthalate (CPET). A typical meal replacement tray and the structure of CPET are shown below:

\textbf{[Diagram omitted: CPET structure]}

Draw the structure of one of the monomers used to make CPET.

\vspace{4cm}

\hfill [1]

\end{parts}

\vfill
\hfill [Total: 4]

\newpage

% ======================================================================
\question A2

Modern methods of chemical analysis often rely on the interpretation of data gathered from instrumental techniques.

Electrophoresis and paper chromatography can both be used to separate amino acids.

Electrophoresis is a separation technique that is based on the mobility of ions in an electric field. Positively charged ions migrate towards a negative electrode and negatively charged ions migrate toward a positive electrode. Ions have different migration rates depending on their total charge, size, and shape, and can therefore be separated.

\textbf{[Diagram omitted: electrophoresis and paper chromatography apparatus]}

\begin{parts}
\part[1] Describe the different principles behind each of the separation techniques.
\vspace{2cm}
\dotfill
\vspace{1cm}
\dotfill [1]

\part[1] Amino acids are colourless. How are the positions of the different amino acids made visible so that measurements for chromatography can be made?
\vspace{1em}
\dotfill [1]

\part[3] This diagram shows the results of a two-way paper chromatography of a mixture of amino acids. The chromatography was carried out using solvent 1 first, then rotated 90\degree{} anticlockwise and carried out for a second time using solvent 2.

\textbf{[Diagram omitted: two-way paper chromatography results]}

\vspace{1em}

\begin{subparts}
\subpart[1] Put a U next to the amino acid with the highest \ce{R_f} value in solvent 2.
\subpart[1] Circle the two amino acids that were not separated in solvent 1.
\subpart[1] Put a W next to the amino acid that was the most soluble in solvent 1.
\end{subparts}

\end{parts}

\vfill
\hfill [Total: 5]

\newpage

% ======================================================================
\question A3

\begin{parts}
\part[2]
\begin{subparts}
\subpart[1] The mass spectrum of the element magnesium is shown below.

\textbf{[Diagram omitted: mass spectrum of magnesium showing peaks at \ce{m/z} 24, 25, 26]}

\vspace{1em}

From the mass spectrum, complete the table with the relative abundances of the three isotopes of magnesium.

\begin{tabular}{|c|c|}
\hline
isotope & relative abundance \\
\hline
\ce{^{24}Mg} & \\
\hline
\ce{^{25}Mg} & \\
\hline
\ce{^{26}Mg} & \\
\hline
\end{tabular}
[1]

\subpart[1] Use your values in (a)(i) to calculate the relative atomic mass, \ce{A_r}, of magnesium.
\vspace{1.5cm}
\dotfill [1]
\end{subparts}

\part[4]
\begin{subparts}
\subpart[2] Similar to Group I carbonates, carbonates of Group II such as magnesium carbonate undergo thermal decomposition to give the metal oxide and carbon dioxide gas.

Predict and explain the trend in thermal stabilities of the carbonates of the Group II elements down the group.
\vspace{3cm}
\dotfill
\vspace{0.5cm}
\dotfill [2]

\subpart[1] When magnesium nitrate, \ce{Mg(NO3)2}, is heated, it readily decomposes and gives off both an acidic gas which dissolves in water to form acid rain as well as a gas that relights a glowing splint.

Write a balanced chemical equation for the action of heat on magnesium nitrate.
\vspace{1cm}
\dotfill [1]

\subpart[1] Suggest a reason why Group I nitrates are unlikely to decompose when heated.
\vspace{1cm}
\dotfill [1]
\end{subparts}

\part[3] X, Y and Z are unknown metals.

A student wants to rank the reactivity of the three metals and has set up a simple cell diagram using the three unknowns.

\textbf{[Diagram omitted: simple cell with metal X, metal Y, and electrolyte \ce{M(NO3)2(aq)}]}

The student made three observations.
(1) Metal Y became smaller after the reaction.
(2) Reddish-brown deposit was found on metal X.
(3) Electrolyte became lighter blue.

\begin{subparts}
\subpart[1] Indicate the flow of electrons on Fig.~1 above.
\vspace{1cm}
\dotfill [1]

\subpart[1] The student concluded that Y was the most reactive metal while Z was the least reactive metal.

Explain why the student's conclusion was wrong.
\vspace{2cm}
\dotfill
\vspace{0.5cm}
\dotfill [1]

\subpart[1] Describe one modification to improve this experiment so that the student could rank the reactivity of the three metals.
\vspace{2cm}
\dotfill
\vspace{0.5cm}
\dotfill [1]
\end{subparts}
\end{parts}

\vfill
\hfill [Total: 9]

\newpage

% ======================================================================
\question A4

Naphtha is a fraction obtained from petroleum (crude oil).

\begin{parts}
\part[3] A student wrote some notes about naphtha.
(1) Naphtha is a pure substance.
(2) Naphtha has a higher boiling point than kerosene.
(3) Naphtha burns with a sootier flame than an alkene with similar carbon chain length.

Explain why his notes were wrong.
\vspace{4cm}
\dotfill
\vspace{0.5cm}
\dotfill
\vspace{0.5cm}
\dotfill [3]

\part[3]
\begin{subparts}
\subpart[1] One compound found in fossil fuel has the formula \ce{C12H26}.

Explain which homologous series this compound belongs to.
\vspace{1.5cm}
\dotfill [1]

\subpart[1] The cracking of \ce{C12H26} is carried out to form an alkene that has six carbon atoms per molecule as one of its products.

Construct a balanced chemical equation for this reaction.
\vspace{1.5cm}
\dotfill [1]
\end{subparts}

\part[2] Ethene can be made by the cracking of hydrocarbons.

Draw a `dot-and-cross' diagram for ethene. You only need to show the outer shell electrons.

[Diagram space]

\vspace{4cm}

\hfill [2]
\end{parts}

\vfill
\hfill [Total: 7]

\newpage

% ======================================================================
\question A5

Molybdenum is a transition element. It is used to make steel that is extremely hard.

Molybdenum can be manufactured by heating together molybdenum(VI) oxide, \ce{MoO3}, and aluminium.

\ce{MoO3 + 2Al -> Al2O3 + Mo}

\begin{parts}
\part[1] Name the reaction between molybdenum(VI) oxide and aluminium.
\vspace{1cm}
\dotfill [1]

\part[2] The process of manufacturing molybdenum is an endothermic reaction. Complete the energy profile diagram below.

\textbf{[Diagram omitted: energy profile diagram axes with reactants labelled]}

Your diagram should include:
\begin{itemize}
\item the names of the products of the reaction,
\item labels to show the enthalpy change of reaction and the activation energy.
\end{itemize}

\vspace{4cm}

\hfill [2]

\part[2] Using ideas about oxidation states, explain why this is a redox reaction.
\vspace{3cm}
\dotfill
\vspace{0.5cm}
\dotfill [2]

\part[1] What is the mass of molybdenum that can be made from \qty{125}{\g} of molybdenum(VI) oxide, \ce{MoO3}?
\vspace{2cm}
\dotfill
\vspace{0.5cm}
\dotfill [1]

\part[1] Predict if the manufacture of molybdenum can be replaced by heating molybdenum(VI) oxide, \ce{MoO3}, with copper.

Explain your answer.
\vspace{1.5cm}
\dotfill
\vspace{0.5cm}
\dotfill [1]

\part[3] Molybdenum has a melting point of \qty{2623}{\celsius}.

With the aid of a labelled diagram, describe the structure and bonding in molybdenum which explains its high melting point.

[Diagram space]

\vspace{5cm}

\hfill [3]
\end{parts}

\vfill
\hfill [Total: 10]

\newpage

% ======================================================================
\question A6

Peroxodisulfate ions, \ce{S2O8^{2-}}, react with iodide ions in aqueous solution. Iron(III) ions, \ce{Fe^{3+}}, catalyse the reaction.

\ce{S2O8^{2-}(aq) + 2I-(aq) -> 2SO4^{2-}(aq) + I2(aq)}

The table shows how the relative rate of the reaction changes when different concentrations of peroxodisulfate ions and iodide ions are used.

\begin{tabular}{|c|c|c|c|c|}
\hline
experiment & concentration of \ce{I-} in \qty{}{\mol\per\dm\cubed} & concentration of \ce{S2O8^{2-}} in \qty{}{\mol\per\dm\cubed} & concentration of \ce{Fe^{3+}} in \qty{}{\mol\per\dm\cubed} & relative rate of reaction \\
\hline
A & 0.01 & 0.004 & 0.01 & 1.7 \\
B & 0.01 & 0.008 & 0.01 & 3.3 \\
C & 0.01 & 0.016 & 0.01 & 6.8 \\
D & 0.02 & 0.004 & 0.02 & 3.4 \\
E & 0.04 & 0.004 & 0.04 & 6.9 \\
F & 0.04 & 0.004 & 0.02 & 6.9 \\
\hline
\end{tabular}

\begin{parts}
\part[2] Use the information in the table to describe how decreasing the concentration of iodide ions affects the relative rate of reaction.
\vspace{3cm}
\dotfill
\vspace{0.5cm}
\dotfill [2]

\part[2]
\begin{subparts}
\subpart[1] Using information from experiment E and F, suggest a reason why the relative rate of the reaction is the same, even though different concentrations of \ce{Fe^{3+}} ion were used.
\vspace{1.5cm}
\dotfill [1]

\subpart[1] Explain how \ce{Fe^{3+}} catalyst affects the rate of the reaction.
\vspace{1.5cm}
\dotfill [1]
\end{subparts}

\part[2] Iron(III) ions can be formed by the reaction of bromine with iron(II) ions.
\ce{Br2(aq) + 2Fe^{2+}(aq) -> 2Br-(aq) + 2Fe^{3+}(aq)}

\begin{subparts}
\subpart[1] Explain the role of bromine in this reaction.
\vspace{1.5cm}
\dotfill [1]

\subpart[1] Using ideas about the strength of oxidising agents, predict if chlorine would oxidise iron(II) ions in a similar manner. Explain.
\vspace{2cm}
\dotfill
\vspace{0.5cm}
\dotfill [1]
\end{subparts}

\part[4] The table shows some iron salts and their uses.

\begin{tabular}{|c|c|}
\hline
salt & use \\
\hline
iron(II) sulfate & iron supplement \\
iron(III) sulfate & dye \\
\hline
\end{tabular}

With the aid of two ionic equations, explain how the two salts in the table can be distinguished from one another.
\vspace{4cm}
\dotfill
\vspace{0.5cm}
\dotfill
\vspace{0.5cm}
\dotfill
\vspace{0.5cm}
\dotfill [4]
\end{parts}

\vfill
\hfill [Total: 10]

\newpage

% ======================================================================
\question A7

Nitrogen oxides in the upper atmosphere cause damage to the ozone layer. Aircraft engines are one source of nitrogen oxides.

\begin{parts}
\part[2] Explain how nitrogen oxides are formed in the engines of an aircraft.
\vspace{2cm}
\dotfill
\vspace{0.5cm}
\dotfill [2]

\part[3]
\begin{subparts}
\subpart[1] Nitrogen oxides are removed from car exhaust emissions by catalytic converters.

Nitrogen oxides can react with unburnt hydrocarbons to form ozone in the presence of sunlight. Briefly state the effect of ozone on human health.
\vspace{1.5cm}
\dotfill [1]

\subpart[1] Write a balanced chemical equation for the simultaneous removal of the harmful pollutants in catalytic converters.
\vspace{1.5cm}
\dotfill [1]

\subpart[1] At high temperatures, nitrogen monoxide reacts with hydrogen according to the following equation.

\ce{2H2(g) + 2NO(g) -> 2H2O(g) + N2(g)}

Suggest a reason why the catalytic converter does not make use of this reaction to remove nitrogen monoxide from car exhaust emissions.
\vspace{1.5cm}
\dotfill [1]
\end{subparts}

\end{parts}

\vfill
\hfill [Total: 5]

\end{questions}

\vfill
\begin{center}
\textbf{-- End of Section A --}
\end{center}

\newpage

% ======================================================================
% SECTION B
% ======================================================================

\section*{Section B}

Answer \textbf{all} three questions in this section. The last question is in the form of an either/or and only one of the alternatives should be attempted.

\begin{questions}

% ======================================================================
\setcounter{question}{0}
\question B8 \textbf{Hess's Law}

Hess's Law states that the overall enthalpy change for a reaction is the same, whether the reaction is carried out in one step or several steps.

This means that the standard enthalpy of reaction depends only on the difference between the standard enthalpy change of the reactants and the standard enthalpy change of the products. It does not depend on the route by which the reaction occurs.

In the enthalpy diagram in Fig.~1, the standard enthalpy change in going from A to B by the direct route is $\Delta H_1$. The standard enthalpy change in going from A to C to B by the indirect route is the sum of $\Delta H_2$ and $\Delta H_3$.

Since the standard enthalpy change is the same for the different routes, this means that
\[\Delta H_1 = \Delta H_2 + \Delta H_3\]

\textbf{[Diagram omitted: enthalpy diagram showing A $\rightarrow$ B with $\Delta H_1$, and A $\rightarrow$ C $\rightarrow$ B with $\Delta H_2$, $\Delta H_3$]}

\textbf{Enthalpy of solution of anhydrous salts}

The enthalpy of solution is the enthalpy change when one mole of ionic substance (e.g. \ce{NaCl}) dissolves to give a dilute solution.

\ce{NaCl(s) -> Na+(aq) + Cl-(aq)}

Magnesium sulfate can exist in either the hydrated form or in the anhydrous form.

Two students investigated the enthalpy of solution of anhydrous magnesium sulfate, $\Delta H_{\text{s}}$, by measuring the initial and highest temperature reached when anhydrous magnesium sulfate was dissolved in water. They presented their results in the following table.

\begin{tabular}{|c|c|}
\hline
mass of anhydrous magnesium sulfate / \qty{}{\g} & 3.01 \\
\hline
volume of water / \qty{}{\cm\cubed} & 50.0 \\
\hline
initial temperature / \qty{}{\celsius} & 17.0 \\
\hline
highest temperature reached / \qty{}{\celsius} & 26.7 \\
\hline
\end{tabular}

\begin{parts}

\part[2] Using the formula provided, calculate the enthalpy change, $\Delta H_{\text{s}}$ in \qty{}{\kJ\per\mol}, when anhydrous magnesium sulfate dissolves in water.

\[\text{enthalpy change (kJ/mol)} = \frac{\text{volume of water (cm}^3) \times 0.0042 \times \text{temperature rise (\celsius)}}{\text{number of moles of anhydrous magnesium sulfate (mol)}}\]

\vspace{4cm}
\dotfill
\vspace{0.5cm}
\dotfill [2]

\part The students repeated the experiment using \qty{6.16}{\g} of solid hydrated magnesium sulfate, \ce{MgSO4*7H2O} and \qty{50.0}{\cm\cubed} of water. The enthalpy change, $\Delta H_{\text{b}}$, was found to be \qty{+18}{\kJ\per\mol}.

The enthalpy change of solution of solid anhydrous magnesium sulfate is difficult to determine experimentally but can be calculated using the enthalpy diagram below.

\textbf{[Diagram omitted: enthalpy cycle showing \ce{MgSO4*7H2O(s)} $\xrightarrow{\Delta H_{\text{b}}}$ \ce{Mg^{2+}(aq) + SO4^{2-}(aq)}, \ce{MgSO4(s) + 7H2O(l)} $\xrightarrow{\Delta H_{\text{s}}}$ same products, and \ce{MgSO4(s) + 7H2O(l)} $\xrightarrow{\Delta H}$ \ce{MgSO4*7H2O(s)}]}

\begin{subparts}
\subpart[1] Calculate the enthalpy change, $\Delta H$, in \qty{}{\kJ\per\mol}, of solid magnesium sulfate, \ce{MgSO4}.
\vspace{2cm}
\dotfill [1]

\subpart[2] The literature value for the enthalpy of solution of anhydrous magnesium sulfate is \qty{-103}{\kJ\per\mol}.

Suggest two reasons for the difference with the value calculated in (b)(i).
\vspace{3cm}
\dotfill
\vspace{0.5cm}
\dotfill [2]
\end{subparts}

\part[2]
\begin{subparts}
\subpart[1] The enthalpy of hydration is the enthalpy change when one mole of gaseous ions dissolve to give a dilute solution.

For example, the enthalpy of hydration of sodium and chloride ions can be represented by the combined equation below:

\ce{Na+(g) + Cl-(g) -> Na+(aq) + Cl-(aq)}

Based on this definition, write an equation, with state symbols, to represent the enthalpy of hydration of magnesium and sulfate ions.
\vspace{1.5cm}
\dotfill [1]

\subpart[1] Enthalpy of hydration is always negative.

Use your equation in (c)(i) to explain why it is always exothermic.
\vspace{2cm}
\dotfill
\vspace{0.5cm}
\dotfill [1]
\end{subparts}

\part[5]
\begin{subparts}
\subpart[3] Magnesium sulfate, commonly known as Epsom salt, is a natural remedy for a number of ailments. Epsom salt has numerous health benefits.

Describe how a pure and dry sample of hydrated magnesium sulfate could be prepared, using magnesium carbonate as one of the starting reagents.
\vspace{4cm}
\dotfill
\vspace{0.5cm}
\dotfill
\vspace{0.5cm}
\dotfill [3]

\subpart[2] In another experiment, \qty{0.100}{\g} of magnesium ribbon is added to \qty{50.0}{\cm\cubed} of \qty{4.00}{\mol\per\dm\cubed} sulfuric acid.

The graph shows the volume of hydrogen produced against time under these experimental conditions.

\textbf{[Diagram omitted: graph of volume of hydrogen versus time]}

Sketch two curves labelled I and II (on the graph above), to show how the volume of hydrogen produced (under the same temperature and pressure) changes with time when:
\begin{itemize}
\item I -- \qty{0.100}{\g} magnesium powder is used instead of the magnesium ribbon;
\item II -- \qty{0.100}{\g} of magnesium ribbon is added to \qty{50}{\cm\cubed} of \qty{0.500}{\mol\per\dm\cubed} sulfuric acid.
\end{itemize}
\hfill [2]
\end{subparts}

\end{parts}

\vfill
\hfill [Total: 12]

\newpage

% ======================================================================
\question B9

The structure of graphite and boron nitride are shown below.

\textbf{[Diagram omitted: graphite structure with key showing carbon atoms; boron nitride structure with key showing boron and nitrogen atoms]}

\begin{parts}
\part[2] Like graphite, boron nitride feels slippery to the touch.

Explain, in terms of structure and bonding, why boron nitride feels slippery to the touch.
\vspace{3cm}
\dotfill
\vspace{0.5cm}
\dotfill [2]

\part[6] Dilute sulfuric acid can be electrolysed using graphite electrodes.

\textbf{[Diagram omitted: electrolysis apparatus with dilute sulfuric acid and graphite electrodes, showing Gas A and Gas B]}

\begin{subparts}
\subpart[1] Graphite is a good electrical conductor. Explain why.
\vspace{2cm}
\dotfill
\vspace{0.5cm}
\dotfill [1]

\subpart[2] Construct the ionic half equations for the reactions happening at the electrodes.

anode:
\vspace{1cm}
\dotfill

cathode:
\vspace{1cm}
\dotfill [2]

\subpart[1] Explain why the volume of gas B produced is approximately double that of gas A.
\vspace{2cm}
\dotfill
\vspace{0.5cm}
\dotfill [1]

\subpart[2] The same set-up can be used for the electrolysis of concentrated hydrochloric acid.

Give one similarity and one difference between the products of the electrolysis of dilute sulfuric acid and concentrated hydrochloric acid.

similarity:
\vspace{1cm}
\dotfill

difference:
\vspace{1cm}
\dotfill [2]
\end{subparts}
\end{parts}

\vfill
\hfill [Total: 8]

\newpage

% ======================================================================
\question B10 \textbf{EITHER}

2-methylpropane (\ce{M_r} = 58) is a hydrocarbon.

\textbf{[Diagram omitted: structure of 2-methylpropane]}

\begin{parts}
\part[4]
\begin{subparts}
\subpart[1] 2-methylpropane reacts with chlorine in the presence of ultraviolet light to give several compounds.

One of these compounds has a relative molecular mass of 127.

Suggest and draw the full structural formula for this compound.
\vspace{3cm}
\dotfill [1]

\subpart[2] 2-methylpropane undergoes complete combustion according to the equation:
\ce{2C4H10(g) + 13O2(g) -> 8CO2(g) + 10H2O(l)}

Calculate the volume of gas produced measured at room temperature and pressure if \qty{1.00}{\kg} of 2-methylpropane was burnt.
\vspace{3cm}
\dotfill
\vspace{0.5cm}
\dotfill [2]

\subpart[1] A single organic compound is formed when one molecule of 2-methylpropane reacts with ten molecules of chlorine.

Write a balanced chemical equation for this reaction.
\vspace{1.5cm}
\dotfill [1]
\end{subparts}

\part[6] The structures of propene and cyclopropane are shown.

\textbf{[Diagram omitted: structures of propene (\ce{CH3CH=CH2}) and cyclopropane (cyclic \ce{C3H6})]}

\begin{subparts}
\subpart[1] Explain why these two compounds are structural isomers of each other.
\vspace{1.5cm}
\dotfill [1]

\subpart[2] Draw the structural formula of all the possible products when propene is reacted with steam under suitable conditions.
\vspace{3cm}
\dotfill [2]

\subpart[1] A student commented that hydrogen could be added to propene to obtain cyclopropane.

Do you agree with the student? Explain why.
\vspace{2cm}
\dotfill
\vspace{0.5cm}
\dotfill [1]

\subpart[2] Describe a chemical test to distinguish between propene and cyclopropane.

test:
\vspace{1cm}
\dotfill

results:
\vspace{1cm}
\dotfill [2]
\end{subparts}
\end{parts}

\vfill
\hfill [Total: 10]

\vfill
\begin{center}
\textbf{-- End of Section B (EITHER Option) --}
\end{center}

\newpage

% ======================================================================
\question B10 \textbf{OR}

Compound A has the structure below.

\textbf{[Diagram omitted: structure of ethyl ethanoate/ester: \ce{CH3COOCH2CH3}]}

\begin{parts}
\part[1] Name compound A.
\vspace{1cm}
\dotfill [1]

\part[3] Compound A reacts with hot aqueous sodium hydroxide to give two compounds, B and C.

\begin{subparts}
\subpart[2] Compound B has the percentage composition by mass:
29.3\% carbon, 3.7\% hydrogen, 39.0\% oxygen, 28.0\% sodium.

Calculate the empirical formula of this compound.
\vspace{3cm}
\dotfill
\vspace{0.5cm}
\dotfill [2]

\subpart[1] Compound C has a relative molecular mass of 74 and is oxidised by acidified potassium manganate(VII) on warming to form butanoic acid.

Suggest and draw the full structural formula for compound C.
\vspace{3cm}
\dotfill [1]
\end{subparts}

\part[2] A student attempted to use an alkane as the starting material to make either an alcohol or a carboxylic acid via a series of two or three step reactions respectively.

\begin{subparts}
\subpart[1] Fill in the blanks below to show how ethanol can be synthesised using a series of two step reactions, using ethane as the starting material.

\begin{tabular}{|l|l|}
\hline
reaction 1: & dehydrogenation (removal of 2 hydrogen atoms) \\
\hline
name of reaction 2: & \dotfill \\
\hline
name of intermediate product: & \dotfill \\
\hline
 & $\rightarrow$ ethanol \\
\hline
\end{tabular}
\hfill [1]

\subpart[1] Ethanol can also be synthesised using fermentation.

State the reagents and conditions required for fermentation to occur.

reagent:
\vspace{1cm}
\dotfill

conditions:
\vspace{1cm}
\dotfill [1]
\end{subparts}

\part[4]
\begin{subparts}
\subpart[1] Explain why compound A is a saturated compound.
\vspace{1.5cm}
\dotfill [1]

\subpart[1] Compound A is used in perfumes.

With reference to a physical property of compound A, explain why it is used in perfumes.
\vspace{1.5cm}
\dotfill [1]

\subpart[1] Draw an isomer of compound A with the same functional group.
\vspace{2cm}
\dotfill [1]

\subpart[1] Suggest one other use for compound A.
\vspace{1cm}
\dotfill [1]
\end{subparts}
\end{parts}

\vfill
\hfill [Total: 10]

\end{questions}

\vfill
\begin{center}
\textbf{-- End of Section B --}
\end{center}

\newpage

% ======================================================================
% ANSWER KEY
% ======================================================================

\section*{Answer Key}

\subsection*{Paper 1 -- Multiple Choice}

\begin{tabular}{|c|c|c|c|c|c|c|c|c|c|c|}
\hline
Q & Ans & Q & Ans & Q & Ans & Q & Ans & Q & Ans \\
\hline
1 & C & 9 & B & 17 & C & 25 & C & 33 & D \\
2 & D & 10 & C & 18 & B & 26 & B & 34 & D \\
3 & A & 11 & B & 19 & D & 27 & C & 35 & B \\
4 & D & 12 & C & 20 & A & 28 & A & 36 & D \\
5 & C & 13 & C & 21 & D & 29 & D & 37 & C \\
6 & D & 14 & A & 22 & D & 30 & C & 38 & D \\
7 & D & 15 & D & 23 & C & 31 & A & 39 & D \\
8 & C & 16 & B & 24 & C & 32 & D & 40 & A \\
\hline
\end{tabular}

\subsection*{Section A}

\textbf{A1}
\begin{itemize}
\item[(a)] B [1]
\item[(b)] Addition -- E \hfill AND \hfill polyester -- C (both correct to gain [1])
\item[(c)] A [1]
\item[(d)]

\textbf{[Diagram omitted: either benzene-1,4-dicarboxylic acid (terephthalic acid) or ethane-1,2-diol (ethylene glycol)]}

\hfill either 1 monomer to gain [1]
\end{itemize}

\textbf{A2}
\begin{itemize}
\item[(a)] Electrophoresis causes separation due to different migration rate in an electric field while chromatography causes separation due to difference in solubility in a solvent. [1]
\item[(b)] Use of a locating agent [1]
\item[(c)] \textbf{[Diagram omitted: annotated chromatography results with U, circles, W]}
\end{itemize}

\textbf{A3}
\begin{itemize}
\item[(a)(i)] (total abundance must add up to 100\%)
\begin{tabular}{|c|c|}
\hline
isotope & relative abundance \\
\hline
\ce{^{24}Mg} & 78 \\
\hline
\ce{^{25}Mg} & 10 \\
\hline
\ce{^{26}Mg} & 12 \\
\hline
\end{tabular}
[1]

\item[(a)(ii)] $\frac{78 \times 24}{100} + \frac{10 \times 25}{100} + \frac{12 \times 26}{100} = 24.34 = 24.3$ [1]

\item[(b)(i)] Thermal stability of Group II carbonates will increase down the group. Down the group, metals increase in reactivity [1], so they will form carbonates that are more stable [1].

\item[(b)(ii)] \ce{2Mg(NO3)2 -> 2MgO + 4NO2 + O2} [1]

\item[(b)(iii)] Group I metals are very reactive and so Group I nitrates are more thermally stable than Group II nitrates. [1]

\item[(c)(i)] Electron flows from Y to X. [1]

\item[(c)(ii)] Electrolyte does not contain ion of Z, so there is no selective discharge, so we are unable to tell whether X or Z is more reactive. OR When Z ion is discharged at electrode X, it only shows that Z is less reactive than H but does not show that it is less reactive than X. [1]

\item[(c)(iii)] Use an electrolyte that contains both ions of X and Z OR remove the voltmeter and connecting wires (to turn it into a displacement reaction). [1]
\end{itemize}

\textbf{A4}
\begin{itemize}
\item[(a)] Naphtha contains a mixture of hydrocarbons / is a mixture of alkanes / a mixture of hydrocarbons or alkanes with similar carbon-chain length or size [1]. Naphtha contains hydrocarbons that are smaller / lighter / shorter than kerosene, so their boiling point should be lower [1]. Naphtha consists of alkanes only; alkanes will have a smaller percentage of carbon compared to alkenes with a similar carbon-chain length, so easier for complete combustion to occur, hence will burn with a less sooty flame [1].

\item[(b)(i)] Alkane, because it follows the general formula \ce{C_nH_{2n+2}} [1]

\item[(b)(ii)] \ce{C12H26 -> C6H12 + C6H14} or \ce{C12H26 -> 2C6H12 + H2} [1]

\item[(c)] \textbf{[Diagram omitted: dot-and-cross diagram of ethene showing outer shell electrons] \hfill [1 mark for showing electrons between C and C; 1 mark for showing electrons between C and H]}
\end{itemize}

\textbf{A5}
\begin{itemize}
\item[(a)] Reduction / Reduction by a more reactive metal [1]

\item[(b)] \textbf{[Diagram omitted: energy profile diagram with reactants, products, \ce{Al2O3 + Mo}, activation energy (Ea) label, enthalpy change ($\Delta H$) label, upward arrow for $\Delta H$]}
\hfill [1 for showing energy level of product and upward arrow; 1 for drawing and labelling arrow for Ea]

\item[(c)] Mo is reduced because its oxidation state decreases from +6 in \ce{MoO3} to 0 in Mo [1]. Al is oxidised because its oxidation state increases from 0 in Al to +3 in \ce{Al2O3} [1]. Since oxidation and reduction occurred simultaneously, the reaction is a redox reaction.

\item[(d)] Mass of Mo $= \frac{96}{144} \times 125\,\text{g} = 83.3\,\text{g}$ [1] \hfill OR \hfill mol of \ce{MoO3} $= 125/144 = 0.8681$ mol; mol of Mo $= 0.8681$; mass of Mo $= 0.8681 \times 96 = 83.3$ g [1]

\item[(e)] No. Copper is a very unreactive metal; it is unlikely to be able to reduce molybdenum(VI) oxide to become molybdenum. [1]

\item[(f)] \textbf{[Diagram omitted: labelled diagram of giant metallic lattice structure of molybdenum]}
\hfill [1 for diagram]
It has a giant metallic lattice structure [1] where strong electrostatic forces of attraction exist between the metal cations held in the lattice and the sea of delocalised and mobile electrons. Hence, a lot of energy is required to overcome the strong electrostatic forces of attraction [1].
\end{itemize}

\textbf{A6}
\begin{itemize}
\item[(a)] Halving the concentration of iodide ions from 0.02 to 0.01 halves the relative rate of reaction from 3.4 to 1.7 [1]. Students will need to show data to gain the second mark (can quote expt A and D or D and E) [1].

\item[(b)(i)] Catalysts will increase rate of reactions when used in small amounts / quantities. [1] (allow for alternative phrasing)

\item[(b)(ii)] \ce{Fe^{3+}} increases rate of a reaction by providing an alternative pathway of lower activation energy without itself being chemically changed. [1]

\item[(c)(i)] It is an oxidising agent. It causes \ce{Fe^{2+}} to gain oxidation state from +2 in \ce{Fe^{2+}} to +3 in \ce{Fe^{3+}} or it causes \ce{Fe^{2+}} to lose electron to become \ce{Fe^{3+}}. [1]

\item[(c)(ii)] Yes. Chlorine is higher up the group than bromine / chlorine is more reactive than bromine, so it is a better / stronger oxidising agent than bromine. [1]

\item[(d)] Sodium hydroxide will produce an insoluble green / dirty green precipitate with iron(II) solution [1] and will produce an insoluble reddish-brown / brown precipitate with iron(III) solution [1].

\ce{Fe^{2+}(aq) + 2OH-(aq) -> Fe(OH)2(s)} [1]

\ce{Fe^{3+}(aq) + 3OH-(aq) -> Fe(OH)3(s)} [1]
\end{itemize}

\textbf{A7}
\begin{itemize}
\item[(a)] Under high temperature in the engine [1], nitrogen and oxygen from air [1] will combine to form nitrogen oxides.

\item[(b)(i)] It can irritate the eyes and lungs OR can cause respiratory problems and lung damage. [1]

\item[(b)(ii)] \ce{2NO + 2CO -> N2 + 2CO2} [1]

\item[(b)(iii)] It is not economical to use hydrogen in this manner / hydrogen is very expensive to obtain / hydrogen is a limited resource / hydrogen is explosive. [1]
\end{itemize}

\subsection*{Section B}

\textbf{B8}
\begin{itemize}
\item[(a)] no. of moles $= 3.01/120 = 0.02508 = 0.0251$ mol; temp rise $= 9.7$ \celsius{} [1]; enthalpy change, $\Delta H_{\text{s}} = -(50 \times 4.2 \times 9.7)/0.02508 = \qty{-81.2}{\kJ\per\mol}$ [1]

\item[(b)(i)] $\Delta H + \Delta H_{\text{b}} = \Delta H_{\text{s}}$; $\Delta H = -81.2 - 18 = \qty{-99.2}{\kJ\per\mol}$ [1]

\item[(b)(ii)] Accept any 2 reasons below [1 mark for 1 reason, max 2 marks]:
\begin{itemize}
\item \ce{MgSO4} not completely anhydrous / OWTTE
\item \ce{MgSO4} is impure
\item heat loss to the atmosphere / surroundings
\item specific heat capacity of solution is taken as that of pure water
\item experiment was done once only so it is not scientific
\item density of solution is taken to be \qty{1}{\g\per\cm\cubed}
\item mass of \ce{7H2O} ignored in calculation
\item uncertainty of thermometer is high so temperature change is unreliable
\item literature values determined under standard conditions but this experiment is not
\item all solid not dissolved
\end{itemize}

\item[(c)(i)] \ce{Mg^{2+}(g) + SO4^{2-}(g) -> Mg^{2+}(aq) + SO4^{2-}(aq)} [1]

\item[(c)(ii)] The particles of magnesium sulfate change from gaseous state to aqueous state, so the particle distance decreases / particles become closer to one another, so heat must be lost for that to happen, hence hydration is exothermic. [1]

\item[(d)(i)] Add excess magnesium carbonate to dilute sulfuric acid and filter to obtain filtrate [1]; heat filtrate to saturation, leave saturated solution to cool and crystallise [1]; filter to obtain magnesium sulfate crystals, rinse with little cold distilled water, dry with filter paper [1].

\item[(d)(ii)] \textbf{[Diagram omitted: I: line which is steeper / increases faster and finishes at the same height; II: line which is less steep / increases more slowly and finishes at the same height] \hfill [1 mark for graph I, 1 mark for graph II] (students need to know that magnesium is the limiting reagent, acid is in excess)}
\end{itemize}

\textbf{B9}
\begin{itemize}
\item[(a)] It has a giant covalent structure [1]. Its atoms are bonded in layers which are held together by weak intermolecular forces which require little amount of energy to overcome [1], hence it is slippery.

\item[(b)(i)] In graphite, there exist mobile electrons [1], where every 1 C atom contributes one valence electron that is not involved in bonding.

\item[(b)(ii)] anode: \ce{4OH-(aq) -> 2H2O(l) + O2(g) + 4e-} [1] \hfill cathode: \ce{2H+(aq) + 2e- -> H2(g)} [1]

\item[(b)(iii)] From the overall equation, we can see that the mole ratio of hydrogen to oxygen is 2:1, so the volume produced will also be in the same ratio. [1]

\item[(b)(iv)] Similarity: Both produce hydrogen gas at the cathode. [1] \hfill Difference: Electrolysis of dilute sulfuric acid produces oxygen at the anode while the other produces chlorine gas at the anode. [1] OR Electrolyte of dilute sulfuric acid causes pH of electrolyte to remain the same while the other causes pH to increase / become more alkaline. [1]
\end{itemize}

\textbf{B10 EITHER}
\begin{itemize}
\item[(a)(i)] \textbf{[Diagram omitted: chlorinated 2-methylpropane with one chlorine substituent, \ce{C4H9Cl}] [1]}

\item[(a)(ii)] Mol of 2-methylpropane $= 1000/58 = 17.24$ mol [1]; mol of carbon dioxide gas $= 68.96$ mol; vol of gas produced $= 68.96 \times 24 = 1655\text{--}1660$ \qty{}{\dm\cubed} (3 s.f.) [1]

\item[(a)(iii)] \ce{C4H10 + 10Cl2 -> C4Cl10 + 10HCl} [1]

\item[(b)(i)] They have the same molecular formula of \ce{C3H6} but different structures / same molecular formula but different arrangement of atoms. [1]

\item[(b)(ii)] \textbf{[Diagram omitted: propan-1-ol (\ce{CH3CH2CH2OH}) and propan-2-ol (\ce{CH3CH(OH)CH3})] \hfill [1 mark for each product] [2]}

\item[(b)(iii)] Disagree. Propene molecular formula is \ce{C3H6}; if hydrogen is added, the product will have molecular formula \ce{C3H8} but cyclopropane has a molecular formula of \ce{C3H6} [1] / product obtained will be propane, not cyclopropane [1].

\item[(b)(iv)] test: Add aqueous bromine / liquid bromine / bromine dissolved in \ce{CCl4} [1]. results: Brown bromine decolourises / brown bromine turns colourless in propene while it remains brown with cyclopropane [1].
\end{itemize}

\textbf{B10 OR}
\begin{itemize}
\item[(a)] Butyl ethanoate [1]

\item[(b)(i)]
\begin{tabular}{|c|c|c|c|c|}
\hline
 & C & H & O & Na \\
\hline
no. of mol & 29.3/12 = 2.44 & 3.7/1 = 3.7 & 39/16 = 2.44 & 28/23 = 1.22 \\
\hline
ratio of mol & 2.44/1.22 = 2 & 3.7/1.22 = 3 & 2.44/1.22 = 2 & 1.22/1.22 = 1 \\
\hline
\end{tabular}

empirical formula = \ce{C2H3O2Na} [1+1=2]

\item[(b)(ii)] \textbf{[Diagram omitted: butan-1-ol, \ce{CH3CH2CH2CH2OH}] [1]}

\item[(c)(i)]
\begin{tabular}{|l|l|}
\hline
reaction 1: & dehydrogenation (removal of 2 hydrogen atoms) \\
\hline
name of reaction 2: & hydration / addition with steam \\
\hline
name of intermediate product: & ethene \\
\hline
 & $\rightarrow$ ethanol \\
\hline
\end{tabular}
\hfill [All 2 correct to gain 1 mark]

\item[(c)(ii)] Reagent: glucose / \ce{C6H12O6} \hfill Conditions: \qty{37}{\celsius}, absence of oxygen, yeast \hfill [Both correct to gain 1 mark]

\item[(d)] It does not contain a carbon-carbon double bond / has only single carbon-carbon bonds. [1]

\item[(e)(i)] It has a low boiling point / it vaporises / evaporates easily at room temperature. [1]

\item[(e)(ii)] any correct isomer, e.g. \ce{CH3CH2CO2CH2CH2CH3}, \ce{CH3CH2CH2CO2CH2CH3}, \ce{CH3CH2CH2CH2CO2CH3}, \ce{HCO2CH2CH2CH2CH2CH3} [1]

\item[(e)(iii)] Solvent / flavourings [1]
\end{itemize}

\vfill
\begin{center}
\textbf{-- End of Answer Key --}
\end{center}

\end{document}